\documentclass[12pt]{article}
\usepackage[utf8]{inputenc}
\usepackage[T1]{fontenc}
\usepackage{hyperref}
\title{Tai N\"ua Language Resources for K-12 Curriculum Development}
\author{Compiled by Claude Sonnet 5 (claude-sonnet-5) for Andr\'as Kornai}
\date{}
\begin{document}
\maketitle
\section{Introduction}
Tai N\"ua (Tai Nuea, Dehong Dai, Tai Mau; ISO 639-3 tdd, Glottocode tain1252) is a Tai-Kadai language of China and Myanmar. It has its own script (Dehong Dai or Tai Le), whose final scheme was released in 1963. This report is based on one web search per language. It found reference and lesson material, but no graded primer or school textbook.
\section{Primary Educational Resources}
No primer or textbook was found in the search results. The SEASite Tai Dehong pages (introduction and lessons, per the Omniglot page) are the closest thing to teaching material:
\begin{itemize}
\item SEASite (Northern Illinois University), Tai Dehong introduction and lessons: \url{https://seasite.niu.edu/tai/TaiDehong/introduction.htm}
\end{itemize}
\section{Linguistic and Reference Materials}
\begin{itemize}
\item Luo Yongxian, \emph{A dictionary of Dehong, Southwest China} (Pacific Linguistics 145): \url{https://search.ub.tum.de/vufind/Record/DE-604.BV013134161}
\item Wikipedia, Tai Nuea language: \url{https://en.wikipedia.org/wiki/Tai_Nuea_language}
\item Omniglot, Tai Nuea script chart and sample text: \url{https://omniglot.com/writing/tainua.htm}
\item Glottolog entry for tdd: \url{https://glottolog.org/resource/languoid/iso/tdd}
\item Unicode proposal on Tai Le script update: \url{https://www.unicode.org/L2/L2024/24098-tai-le-update.pdf}
\end{itemize}
The search also mentioned Yu Cuirong's Tai N\"u--Chinese dictionary (2004), Hudak's edition of Gedney's Southwestern Tai glossaries (1994) and Young's \emph{Shan chrestomathy} (1985). No URLs for these were seen, so none are given.
\section{Religious and Community Resources}
Nothing was found in the search results.
\section{Digital and Audio Resources}
No audio or software resources were found. The Unicode proposal above is relevant to digital text support for the script.
\section{Curriculum Development Context}
The sources show a written standard and a published dictionary. They show no school materials. Curriculum work would have to start from the dictionary, the script documentation and the SEASite lessons. Further searching, especially in Chinese and in Myanmar sources, is needed.
\end{document}